\chapter{Rotary, coolant and remote control}
\label{ch:remote}

Three accessories change how a laser is used rather than what it can burn: a
rotary chuck turns cylinders into engravable objects, coolant and air assist keep
the cut alive, and a network connection lets a job be started from somewhere else.
This chapter covers all three, plus the remote-control servers that make MeerK40t
scriptable from a browser or another machine.

\section{Rotary: engraving on cylinders}

A rotary attachment replaces one axis with a chuck that rotates the workpiece.
The coordination problem is the same on every machine: the circumference of the
object must be mapped onto the axis that used to be the bed's height.

\begin{walkthrough}{Rotary, step by step}
\begin{enumerate}[label=\arabic*.,leftmargin=1.4em]
  \item Open the rotary settings: \menu{Tools>Rotary} (device settings), or the
        \btn{Rotary} button on the device's ribbon.
  \item Tick \btn{Rotary-Mode active}, then --- before anything else --- set
        \btn{Outside diameter (mm)} and \btn{Usable length (mm)}. The fields
        below them depend on these two values.
  \item Leave \btn{Auto Y wrap from diameter} on for your first run. The Y axis
        is then scaled so that the bed height corresponds to the object's
        circumference ($\pi \times \text{diameter}$).
  \item Select your artwork on the scene and click \btn{Fit selection to rotary}.
        It scales the selection uniformly to fit usable length $\times$
        circumference.
  \item Burn a test strip on a piece of scrap tube before committing to the real
        object. Check the join line where the engraving wraps around: if the
        pattern does not meet cleanly, the circumference is wrong --- measure the
        object rather than trusting a stated size.
  \item With rotary mode active, the Home command is suppressed by default
        (\setting{suppress_home}), so the machine does not drive the head to a
        home position that may be occupied by the chuck.
\end{enumerate}
\end{walkthrough}

\begin{center}\small
\begin{tabular}{@{}L{4.6cm}L{1.9cm}L{7.6cm}@{}}
\toprule
\textbf{Setting} & \textbf{Default} & \textbf{Meaning} \\
\midrule
\setting{rotary_active} & off & Master switch for rotary mode. \\
\setting{rotary_diameter_mm} & 80.0 & Outside diameter; $\pi D$ becomes the wrap
height. \\
\setting{rotary_length_mm} & 200.0 & Usable length along the axis. \\
\setting{rotary_auto_wrap_scale} & on & Derive the Y scale from the diameter
rather than using a manual factor. \\
\setting{rotary_auto_length_scale} & off & Derive the X scale from the usable
length. \\
\setting{rotary_scale_x}, \setting{rotary_scale_y} & 1.0 & Manual scales, used
when the automatic ones are off. \\
\setting{rotary_flat_y_steps} & 159.6 & The flat-bed steps/mm to return to
afterwards (GRBL). \\
\setting{rotary_y_steps} & 80.0 & Steps/mm of the rotary axis (GRBL). \\
\setting{rotary_steps_compensate} & on & Compensate in software, leaving the
controller's own settings untouched (GRBL). \\
\setting{rotary_firmware_steps} & off & Write the rotary steps into the
controller at job start and restore them afterwards (GRBL). \\
\setting{rotary_cal_test_mm} & 100.0 & Length of the calibration test burn. \\
\setting{suppress_home} & on & Skip the home command while rotary mode is
active. \\
\setting{rotary_home_x_only} & on & Home only X, never Y, with a chuck fitted
(GRBL). \\
\setting{rotary_flip_x}, \setting{rotary_flip_y} & off & Mirror the axes. \\
\bottomrule
\end{tabular}
\end{center}

\noindent The console side is short and worth knowing, because it is what a
repeatable rotary process ends up being driven by:

\begin{lstlisting}[style=konsole]
rotary                       # show status: active, diameter, circumference, scales
rotaryfit                    # scale emphasised elements to length x circumference
rotaryscale                  # scale by the effective scales
rotarysuggest 80 200 8 1     # estimate steps/mm: diameter, motor steps, microsteps, ratio
rotarycal 314.2              # recalibrate from a measured test burn
\end{lstlisting}

\subsection*{Cylinder correction (galvo machines)}

Galvo markers with a JCZ controller can also engrave curved surfaces without
rotating them, by projecting the design onto the cylinder. This is a
\emph{different} feature, and it applies to Balor devices only:

\begin{lstlisting}[style=konsole]
cylinder axis X 50mm     # the cylinder's diameter, and which axis runs around it
cylinder distance 120mm  # distance from the cylinder to the mirror
cylinder on
\end{lstlisting}

\noindent The matching settings are \setting{cylinder_active},
\setting{cylinder_mirror_distance} (100\unit{mm} by default),
\setting{cylinder_x_axis} / \setting{cylinder_y_axis} and the diameters
\setting{cylinder_x_diameter} / \setting{cylinder_y_diameter} (50\unit{mm}).
Curved-surface correction keeps straight lines straight when the surface falls
away from the lens --- which, for a small tube entirely inside the galvo field,
is often better than rotating it.

\section{Coolant, air assist and sleep}

\subsection*{Coolant and air assist}

Each device has a \btn{Coolant} setting that names how the machine's coolant or
air assist is switched. The built-in methods are:

\begin{itemize}
  \item \emph{Warnmessage} --- shows a dialog asking you to switch the air assist
        on or off at the right moment;
  \item \emph{GCode M7/M9} and \emph{GCode M8/M9} --- for GRBL controllers with a
        controllable output.
\end{itemize}

\noindent Where you use it:

\begin{itemize}
  \item Per operation, with the \emph{Coolant} field on the operation's
        properties: \emph{No changes}, \emph{Turn on} or \emph{Turn off}. Put
        ``turn on'' on your first operation and ``turn off'' on your last.
  \item In the operations tree, with \btn{Append Coolant On} and
        \btn{Append Coolant Off} as explicit steps in the job.
  \item From the ribbon with the \btn{Coolant} button, which is only visible when
        the device has a coolant configured.
\end{itemize}

\noindent The console commands are \cmd{coolants} (list the available methods),
\cmd{coolant\_on}/\cmd{coolant\_off}, and
\cmd{coolant\_on\_by\_id <id>}/\cmd{coolant\_off\_by\_id <id>} for choosing a
method explicitly.

\subsection*{Preventing the computer from sleeping}

A long job that outlives the screen saver is a job that can be interrupted.
\cmd{system\_hibernate prevent} stops the machine from sleeping;
\cmd{system\_hibernate allow} releases it. On Linux this masks the sleep, suspend,
hibernate and hybrid-sleep targets through \cmd{systemctl} --- which is why the
program mentions that administrative privileges may be required --- and on Windows
it uses the standard execution-state API. The device configuration's
\emph{Default Actions} can bind \emph{Hibernation off} and \emph{Hibernation on}
to the start and end of a job.

\section{Remote control over the network}

\menu{Network} starts and stops MeerK40t's own servers. Each one speaks a
different language to a different kind of client:

\begin{center}\small
\begin{tabular}{@{}L{3.2cm}L{3.0cm}L{8.0cm}@{}}
\toprule
\textbf{Server} & \textbf{Default port} & \textbf{What it is for} \\
\midrule
Webserver & 8080 (menu), 2080 (console command default) & A small web console:
device status, the job queue with pause/stop, a command box and the last console
output as HTML. It listens on localhost only. \\
Telnet (\emph{consoleserver}) & 23 & The full MeerK40t console over telnet. Every
command you can type locally, you can type from another machine. \\
GRBL-Server (\emph{grblcontrol}) & 23 & Emulates a GRBL controller, so software
that expects to talk to GRBL (or a sender like a g-code streaming tool) can drive
your machine through MeerK40t. \cmd{grblmock} is a trivial mock version for
testing clients. \\
Ruida-Server (\emph{ruidacontrol}) & UDP 50200, jog 50207 & Emulates a Ruida
controller --- the bridge that lets Ruida-speaking software drive a K40. It can
also work as a man in the middle (40200/40207) or as a LightBurn-to-Ruida bridge
on TCP 5005. \\
\bottomrule
\end{tabular}
\end{center}

\noindent Start them from the menu, or from the console, which is how you script
them:

\begin{lstlisting}[style=konsole]
webserver -p 8080          # start the web console on a chosen port
consoleserver -p 2323      # a second telnet console, no privilege needed
grblcontrol -p 23          # GRBL emulator, spooling to the active device
ruidacontrol               # Ruida emulation (UDP 50200/50207)
\end{lstlisting}

\noindent The web console accepts commands as \cmd{?cmd=<command>} in a GET or
POST request, and job control as \cmd{job\_cmd=<index>:<operation>} where the
operations are \cmd{stop\_job}, \cmd{remove\_job}, \cmd{pause\_device} and
\cmd{resume\_device}. That is enough to build a simple shop-floor dashboard:
pause every machine, see what is running, and start the next job, from a browser.

\begin{warnbox}[Remote control is real control]
A telnet console or an emulated controller accepts commands that move the head,
fire the laser and run jobs. The web server binds to localhost by default for
exactly this reason --- do not expose these ports to a network you do not
control, and never leave a server running on a machine that is also reachable
from the internet.
\end{warnbox}

\subsection*{A networked K40}

Lihuiyu devices can also be driven over the network: in the device's
\tabname{Interface} tab, choose \btn{Networked} and set the \emph{Address} and
\emph{Port} of a machine running MeerK40t as a server (the default ports are 1025
on Linux/macOS and 1022 on Windows). The \btn{Network-Controller} window then
replaces the USB controller window. This is the standard way to put the noisy
machine in one room and the computer in another.

\subsection*{GRBL over Wi-Fi: ESP3D}

Diode machines with an ESP32-based Wi-Fi module --- MKS DLC32 boards are the
common ones --- can be driven without a cable. Enable it in the GRBL
configuration's \tabname{ESP3D Upload} tab: \btn{Enable ESP3D Upload}, then
\btn{ESP3D Host} (default \file{192.168.1.100}), \btn{ESP3D Port} (80),
\btn{ESP3D SD Path} (\file{/}) and optional credentials, with
\btn{Test Connection} and \btn{List Files} showing the SD card's usage. The
ribbon then gains \btn{ESP3D Upload+Run} (right-click to upload without
starting), \btn{ESP3D Pause}, \btn{ESP3D Resume}, \btn{ESP3D Stop} and
\btn{ESP3D Files}. \btn{Prepare Plan exports for SD card} adjusts exported files
for the MKS DLC32 (LF line endings, M3 laser mode, no \cmd{G28}), and
\btn{Cleanup Local Files} deletes the local copy after upload. The console
equivalents are \cmd{esp3d\_config test|set <param> <value>} and
\cmd{esp3d\_list}.

\section{Connecting to other software}

Two bridges are worth knowing about even if you never use them:

\begin{itemize}
  \item \btn{GRBL-Server} lets a phone app or a g-code sender think it is talking
        to a normal GRBL machine while MeerK40t does the work --- useful when you
        want someone else's jogging interface or a CAM program's streamer.
  \item \btn{Ruida-Server} does the same trick in the other direction: software
        that only knows Ruida can drive a K40, because MeerK40t emulates the
        controller the software expects.
\end{itemize}

\begin{tryit}{Drive a job from a browser}
\begin{enumerate}[label=\arabic*.,leftmargin=1.4em]
  \item Start the web server: \menu{Network>Start Webserver}, or
        \cmd{webserver} in the console.
  \item Open a browser at \file{http://127.0.0.1:8080}. You should see the web
        console with the active device, the job list and a command box.
  \item Queue a small job from the desktop interface and watch it appear in the
        web page's job list.
  \item Use the web page to pause and resume the device. The desktop interface
        should follow along --- they are two views of one program.
  \item Stop the server from the \menu{Network} menu, and confirm in the console
        that the port is closed.
  \item If you use a shop network, repeat with \cmd{consoleserver} and a real
        second machine --- then decide, deliberately, whether that convenience is
        worth the risk.
\end{enumerate}
\end{tryit}
